\documentclass{article}
% --- Packages and PDF configuration ---

\usepackage[a5paper,left=12mm,right=12mm,top=12mm,bottom=6mm,includefoot,footskip=8mm]{geometry}
\usepackage{fontspec}
\usepackage{babel}
\babelprovide[import]{russian}
\usepackage{graphicx}
\usepackage{xcolor}
\usepackage{array}
\usepackage{longtable}
\usepackage{colortbl}
\usepackage{fancyhdr}
\usepackage{hyperref}
\hypersetup{
  unicode=true,
  pdfpagelayout=TwoColumnRight,
  colorlinks=false,
  pdfborder={0 0 0}
}

% PDF outline hierarchy matches the unnumbered document sections.
\makeatletter
\def\toclevel@section{0}
\def\toclevel@subsection{1}
\makeatother

% --- Common configuration and styles ---

% Fonts
% Golos Text has no italic face. Use an explicit oblique for emphasis.
\setmainfont{golostext-400.ttf}[Path=fonts/,Ligatures=NoCommon,
  ItalicFont=golostext-400.ttf,ItalicFeatures={FakeSlant}]
\newfontfamily\GolosText{golostext-400.ttf}[Path=fonts/,Ligatures=NoCommon,
  ItalicFont=golostext-400.ttf,ItalicFeatures={FakeSlant}]
\newfontfamily\GolosTextMedium{golostext-500.ttf}[Path=fonts/,Ligatures=NoCommon]
\newfontfamily\JetBrainsMono{jetbrainsmono-400.ttf}[Path=fonts/,Ligatures=NoCommon]
\newfontfamily\TablerDice{tabler-300.ttf}[Path=fonts/]
\urlstyle{same}

% Colors
\definecolor{rulegray}{gray}{0.60}

% Page style
\pagestyle{fancy}
\fancyhf{}
\renewcommand{\headrulewidth}{0pt}
\renewcommand{\footrulewidth}{0pt}
\fancyfoot[C]{{\GolosText\fontsize{9}{11.2}\selectfont\thepage}}

% Paragraph style
\setlength{\parindent}{0pt}
\newlength{\ParagraphSpacing}
\setlength{\ParagraphSpacing}{7.5pt}
\setlength{\parskip}{\ParagraphSpacing}
\newcommand{\ParagraphStyle}{%
  \GolosText\fontsize{9}{12}\selectfont\raggedright
  \setlength{\parindent}{0pt}%
  \setlength{\parskip}{\ParagraphSpacing}}

% Heading spacing
\newlength{\SectionBeforeSpace}
\newlength{\SubsectionBeforeSpace}
\newlength{\SectionAfterSpace}
\newlength{\SubsectionAfterSpace}
\setlength{\SectionBeforeSpace}{12pt}
\setlength{\SubsectionBeforeSpace}{8.5pt}
\setlength{\SectionAfterSpace}{10.5pt}
\setlength{\SubsectionAfterSpace}{\ParagraphSpacing}

% Heading formatting
\newcommand{\HeadingStyle}{\GolosTextMedium\raggedright}
\setcounter{secnumdepth}{0}
\makeatletter
\newcommand{\StartHeading}[6]{%
  % For consecutive headings, use this heading's before-space instead of
  % the previous heading's after-space, which standard sectioning retains.
  \if@nobreak
    \vskip-\lastskip
    \vskip-\dimexpr#4\relax
  \fi
  \@startsection{#1}{#2}{#3}{#4}{#5}{#6}}
\newcommand{\SectionHeading}{%
  \StartHeading{section}{1}{0pt}%
    {-\SectionBeforeSpace}{\SectionAfterSpace}%
    {\HeadingStyle\fontsize{16}{19}\selectfont}}
% Major sections start on a new page. Starred sections do not force a page break.
\renewcommand{\section}{%
  \@ifstar{\SectionHeading*}{\clearpage\SectionHeading}}
\renewcommand{\subsection}{%
  \StartHeading{subsection}{2}{0pt}%
    {-\SubsectionBeforeSpace}{\SubsectionAfterSpace}%
    {\HeadingStyle\fontsize{12}{15}\selectfont}}
\makeatother

% --- Cover ---

\newcommand{\CoverTitle}[1]{%
  {\GolosTextMedium\fontsize{33}{42}\selectfont #1\par}}

\newcommand{\CoverSubtitle}[1]{%
  \vspace{28.5pt}
  {\GolosText\fontsize{13}{16}\selectfont #1\par}}

\newcommand{\CoverEditionLine}[1]{%
  \vfill
  {\GolosText\fontsize{9}{12}\selectfont #1\par}}

\NewDocumentEnvironment{Cover}{+b}{%
  \thispagestyle{empty}
  \AddToHookNext{shipout/foreground}{%
    \put(12mm,-12mm){%
      \begingroup
      \setlength{\fboxsep}{0pt}%
      \setlength{\fboxrule}{0.5pt}%
      \raisebox{-\height}[0pt][0pt]{\color{black}\fbox{%
        \begin{minipage}[t][\dimexpr\paperheight-24mm-2\fboxrule\relax][t]{\dimexpr\paperwidth-24mm-2\fboxrule\relax}
          \centering
          \vspace*{72mm}
          #1%
          \vspace*{10mm}
        \end{minipage}}}%
      \endgroup}}%
  \null
}{}

% --- Publication details ---

\newlength{\ImprintBlockSpace}
\setlength{\ImprintBlockSpace}{\dimexpr\ParagraphSpacing+3pt\relax}

\newenvironment{Imprint}{%
  \thispagestyle{empty}%
  \ParagraphStyle
  \ParagraphAlignment
}{\par}

% Dedicated imprint typography uses paragraph spacing and its own block spacing.
% The following paragraph supplies the gap after titles and labels exactly once.
\newcommand{\ImprintTitle}[1]{%
  {\GolosTextMedium\fontsize{12}{15}\selectfont #1\par}\nobreak}
\newcommand{\ImprintLabel}[1]{%
  \par\addvspace{\ImprintBlockSpace}%
  {\GolosTextMedium #1\par}\nobreak}

\newenvironment{ImprintBlock}{%
  \par\addvspace{\ImprintBlockSpace}%
}{\par}

\newcommand{\ImprintText}[1]{#1\par}

% --- Contents ---

% Entry spacing
\newcommand{\ContentsSubentryIndent}{1.5em}
\newcommand{\ContentsEntryGap}{0.5em}
\newcommand{\ContentsLeaderStep}{0.55em}
\setcounter{tocdepth}{2}

% Entry formatting
\newcommand{\ContentsEntry}[2]{%
  \noindent#1\nobreak\hspace{\ContentsEntryGap}\color{rulegray}%
  \leaders\hbox to \ContentsLeaderStep{.\hss}\hfill\hspace{\ContentsEntryGap}\color{black}#2\par}
\newcommand{\ContentsSubentry}[2]{%
  \ContentsEntry{\hspace*{\ContentsSubentryIndent}#1}{#2}}

% Automatic contents integration
\makeatletter
\renewcommand{\l@section}[2]{\ContentsEntry{#1}{#2}}
\renewcommand{\l@subsection}[2]{\ContentsSubentry{#1}{#2}}

\newcommand{\ContentsEntries}{%
  \begingroup
  \ParagraphStyle
  \@starttoc{toc}%
  \endgroup}
\makeatother

% --- Dictionary description ---

% Paragraph alignment and hyphenation
\newcommand{\ParagraphAlignment}{%
  \selectlanguage{russian}%
  \rightskip=0pt plus 2em
  \parfillskip=0pt plus 1fil
  \spaceskip=\fontdimen2\font
  \xspaceskip=0pt
  \tolerance=1000
  \doublehyphendemerits=5000
  \finalhyphendemerits=5000
  \lefthyphenmin=3
  \righthyphenmin=3\relax}

\newcommand{\TextParagraph}[1]{%
  \begingroup
  \ParagraphStyle
  \ParagraphAlignment
  % Keep each explanatory paragraph on one page without joining adjacent blocks.
  \interlinepenalty=10000\relax
  #1\par
  \endgroup}

% Place reference link targets above the first line so it remains visible after navigation.
\makeatletter
\newcommand{\ReferenceTarget}[1]{\Hy@raisedlink{\hypertarget{#1}{}}}
\makeatother

% Lists
\newlength{\ListIndent}
\newcommand{\NumberedLabelSep}{0.25em}

\makeatletter
\newenvironment{List}{%
  \ParagraphStyle
  \ParagraphAlignment
  % Standard lists restore rightskip from this register during setup.
  \@rightskip=\rightskip
  % Keep lines and items together, leaving the list boundaries breakable.
  \interlinepenalty=10000\relax
  \@itempenalty=10000\relax
  % After a heading, \@nbitem merges its gap with the heading's existing skip.
  % Preserve the paragraph gap separately, as for a following TextParagraph.
  \if@nobreak \vskip\parskip \fi
  \setlength{\ListIndent}{2.5em}%
  \begin{list}{}{%
    % The surrounding paragraph spacing already supplies the external gap.
    \setlength{\topsep}{0pt}%
    \setlength{\partopsep}{0pt}%
    \setlength{\itemsep}{0pt}%
    \setlength{\parsep}{0pt}%
    \setlength{\leftmargin}{\ListIndent}%
    \setlength{\rightmargin}{0pt}%
    \setlength{\itemindent}{0pt}%
    \setlength{\listparindent}{0pt}%
    \setlength{\labelwidth}{\dimexpr\ListIndent-\NumberedLabelSep\relax}}%
}{\end{list}}
\makeatother

\newcommand{\NumberedItem}[2]{%
  \setlength{\labelsep}{\NumberedLabelSep}%
  \item[#1.]#2\par}
\newcommand{\BulletItem}[1]{%
  \setlength{\labelsep}{0.4em}%
  \item[\textbullet]#1\par}

% Standard block spacing uses one shared topsep above and below the table.
\newlength{\TransliterationBlockGap}
\newlength{\TransliterationColumnGap}
\setlength{\TransliterationBlockGap}{7mm}
\setlength{\TransliterationColumnGap}{4mm}

\newenvironment{TransliterationTable}{%
  \ParagraphStyle
  \setlength{\topsep}{\ParagraphSpacing}%
  \setlength{\partopsep}{0pt}%
  \begin{trivlist}%
  \item\relax
  \renewcommand{\arraystretch}{1.2}%
  \noindent\ignorespaces
}{\unskip\end{trivlist}}

\newcommand{\TextTableHeading}[1]{{\GolosTextMedium #1}}

\newenvironment{TransliterationBlock}{%
  \begin{tabular}[t]{@{}l@{\hspace{\TransliterationColumnGap}}l@{}}%
    \TextTableHeading{Буква} & \TextTableHeading{Запись}\\[\TableHeaderAfterSpace]\hline
}{\end{tabular}}

\newcommand{\TransliterationRow}[2]{#1 & #2\\\hline}

% --- Dictionary table ---

% Dimensions and alignment parameters
\setlength{\arrayrulewidth}{0.3pt}
\arrayrulecolor{rulegray}
\newlength{\TableHeaderAfterSpace}
\setlength{\TableHeaderAfterSpace}{0.2mm}
\newlength{\TableColumnPadding}
\setlength{\TableColumnPadding}{1.15mm}
\newlength{\TableRowHeight}
\newlength{\TableRowDepth}
\newlength{\TableDiceVerticalGap}
\newlength{\TableCapitalHeight}
\newlength{\TableTextLift}
\newlength{\TableDiceLeftShift}
\setlength{\TableRowHeight}{5.87mm}
\setlength{\TableRowDepth}{1.5mm}

% Align the visible left contour of the die frame with the left contour of «К».
% The value includes the two fonts' measured side bearings at their output sizes.
\setlength{\TableDiceLeftShift}{-0.198524mm}
\newsavebox{\TableCapitalReferenceBox}
\newsavebox{\TableDiceReferenceBox}

% Reference measurements and vertical alignment
\newcommand{\TableInitializeVerticalAlignment}{%
  \sbox{\TableCapitalReferenceBox}{{\GolosText\fontsize{9}{11.2}\selectfont H}}%
  \setlength{\TableCapitalHeight}{\ht\TableCapitalReferenceBox}%
  \setlength{\TableDiceVerticalGap}{\dimexpr\TableRowHeight/2-2.25mm\relax}%
  \setlength{\TableTextLift}{\dimexpr\TableRowHeight/2-\TableRowDepth-\TableCapitalHeight/2\relax}}

% Dice rendering
\newcommand{\TableDiceSymbol}[1]{%
  \ifnum#1<1
    \PackageError{dictionary}{Dice value #1 is below 1}{}
  \else
    \ifnum#1>6
      \PackageError{dictionary}{Dice value #1 is above 6}{}
    \else
      \char\numexpr"F08A+#1\relax
    \fi
  \fi}

\newcommand{\TableDiceGlyph}[1]{%
  \begingroup
  \sbox{\TableDiceReferenceBox}{\resizebox*{!}{4.5mm}{%
    \TablerDice
    \TableDiceSymbol{#1}}}%
  \raisebox{\dimexpr-\TableRowDepth+\TableDiceVerticalGap+\dp\TableDiceReferenceBox\relax}[0pt][0pt]{\usebox{\TableDiceReferenceBox}}%
  \endgroup}

\newcommand{\TableDiceGlyphBox}[1]{\makebox[5mm][l]{\hspace*{\TableDiceLeftShift}\TableDiceGlyph{#1}}}

\newcommand{\TableDiceHeader}[1]{%
  \TableHeaderFont\hspace*{0.04mm}#1}

% Search and copy support in PDF
\newcommand{\TableDiceUnicodeHex}[1]{%
  \ifcase#1\or 2680\or 2681\or 2682\or 2683\or 2684\or 2685\fi}

\newcommand{\TableDiceActualTextStart}[4]{%
  \begingroup
  \edef\TableDiceActualTextCommand{\endgroup
    \noexpand\pdfextension literal direct
    {/Span << /ActualText <FEFF\TableDiceUnicodeHex{#1}\TableDiceUnicodeHex{#2}\TableDiceUnicodeHex{#3}\TableDiceUnicodeHex{#4}> >> BDC}}%
  \TableDiceActualTextCommand}

\newcommand{\TableDiceSearchText}[1]{%
  \makebox[0pt][l]{%
    \pdfextension literal direct {3 Tr}%
    {\JetBrainsMono\fontsize{9}{11.2}\selectfont #1}%
    \pdfextension literal direct {0 Tr}}}

\newcommand{\TableDice}[1]{\expandafter\TableDiceFour#1\TableDiceEnd}
\def\TableDiceFour#1#2#3#4\TableDiceEnd{%
  \TableDiceSearchText{#1#2#3#4}%
  \TableDiceActualTextStart{#1}{#2}{#3}{#4}%
  \TableDiceGlyphBox{#1}\TableDiceGlyphBox{#2}\TableDiceGlyphBox{#3}\TableDiceGlyphBox{#4}%
  \pdfextension literal direct {EMC}}

% Cells, rows and header
\newcommand{\TableBodyFont}{\GolosText\fontsize{9}{11.2}\selectfont}
\newcommand{\TableHeaderFont}{\GolosText\fontsize{8}{10}\selectfont}
\newcommand{\TableMonospaceFont}{\JetBrainsMono\fontsize{9}{11.2}\selectfont}
\newcommand{\TableRowStrut}{\rule[-\TableRowDepth]{0pt}{\TableRowHeight}}
\newcommand{\TableTextCell}[1]{\raisebox{\TableTextLift}[0pt][0pt]{{\TableBodyFont\mbox{#1}}}}
\newcommand{\TableMonospaceCell}[1]{\raisebox{\TableTextLift}[0pt][0pt]{{\TableMonospaceFont\mbox{#1}}}}

\newcommand{\TableDataRow}[6][]{%
  \TableRowStrut\mbox{\raisebox{\dimexpr\TableRowHeight-\TableRowDepth\relax}[0pt][0pt]{#1}\TableDice{#2}} &
  \TableTextCell{#3} &
  \TableMonospaceCell{#4} &
  \TableMonospaceCell{#5} &
  \TableMonospaceCell{#6}\\\hline}

\newcommand{\TableHeader}[5]{%
  \TableDiceHeader{#1} &
  {\TableHeaderFont #2} &
  {\TableHeaderFont #3} &
  {\TableHeaderFont #4} &
  {\TableHeaderFont #5}\\[\TableHeaderAfterSpace]\hline
  \endhead}

% Table environment and spread positioning
\newcommand{\TableStartOnLeftPage}{%
  \clearpage
  \ifodd\value{page}
    \thispagestyle{empty}\null\clearpage
  \fi}

\newenvironment{DictionaryTable}{%
  \setlength{\tabcolsep}{\TableColumnPadding}%
  \setlength{\LTleft}{0pt}%
  \setlength{\LTright}{0pt}%
  \setlength{\LTpre}{0pt}%
  \setlength{\LTpost}{0pt}%
  \begin{longtable}{@{\extracolsep{\fill}}p{25mm}p{30mm}p{30mm}p{16.5mm}p{10mm}@{}}
}{\end{longtable}}

\input{dictionary-version.tex}
\hypersetup{
  pdftitle={Русский Diceware 4d6},
  pdfauthor={Игорь Малиняк},
  pdfsubject={Словарь для парольных фраз}
}

% --- Document content ---

\begin{document}
  \pdfbookmark[0]{Обложка}{cover}
  \begin{Cover}
    \CoverTitle{Русский\\Diceware 4d6}
    \CoverSubtitle{Словарь для парольных фраз}
    \CoverEditionLine{\DictionaryEdition-е издание \textbullet\ \DictionaryRevision-я редакция}
  \end{Cover}
  \clearpage
  \pdfbookmark[0]{Об издании}{imprint}
  \begin{Imprint}
    \ImprintTitle{Русский Diceware 4d6}
    \ImprintText{Словарь для парольных фраз}
    \begin{ImprintBlock}
      \ImprintText{Автор текста и составитель словаря — Игорь Малиняк}
      \ImprintText{\DictionaryEdition-е издание, \DictionaryRevision-я редакция}
      \ImprintText{2026}
    \end{ImprintBlock}
    \ImprintLabel{Условия использования}
    \ImprintText{Текст и словарь переданы в общественное достояние на условиях CC0 1.0 Universal. Автор отказывается от
      авторских и смежных прав в той мере, в какой это допускает закон.}
    \ImprintText{Текст и словарь можно копировать, изменять и распространять, в том числе в коммерческих целях, без
      запроса разрешения. CC0 не затрагивает права третьих лиц.}
    \ImprintText{Подробнее о CC0: \url{https://creativecommons.org/publicdomain/zero/1.0/}}
    \ImprintText{Шрифты и символы кубиков используются на условиях собственных лицензий.}
    \ImprintLabel{Страница проекта}
    \ImprintText{\url{https://github.com/igor-malinyak/Russian-Diceware-4d6}}
  \end{Imprint}
  \clearpage
  \pdfbookmark[0]{Содержание}{contents}
  \section*{Содержание}
  \ContentsEntries
  \section{Метод Diceware и назначение словаря}
  \TextParagraph{Diceware — это способ создавать парольные фразы с помощью обычных шестигранных кубиков и специального
    списка слов. Каждой комбинации на кубиках соответствует одно слово. Несколько последовательных бросков позволяют
    получить парольную фразу из случайных слов. Метод предложил Арнольд Рейнхольд в 1995 году
    \hyperlink{source-1}{[1]}.}
  \TextParagraph{Стойкость такой фразы основана на случайном выборе слов. При достаточном их количестве фразу трудно
    подобрать, даже если известны словарь и способ создания пароля.}

  \subsection{Язык и объем словаря}
  \TextParagraph{Наиболее известные списки Diceware составлены на английском языке. Для человека, который не является
    его носителем, часть слов может оказаться незнакомой или неудобной для запоминания. Даже знакомые иностранные слова
    бывает труднее связать между собой в яркий образ. Для русскоязычных пользователей словарь на русском языке облегчает
    запоминание фразы благодаря привычным словам и ассоциациям.}
  \TextParagraph{Список для пяти кубиков содержит 7776 слов. При таком объеме сложнее обойтись без редких и специальных
    слов, а также без слов, слишком близких друг к другу по смыслу. В Русском Diceware 4d6 всего 1296 слов — по одному
    на каждую возможную комбинацию четырех кубиков. При отборе предпочтение отдавалось широко известным словам, которые
    легко представить. Однако меньший размер словаря означает, что для сопоставимой защиты от подбора потребуется больше
    слов во фразе.}

  \subsection{Количество слов в пароле}
  \TextParagraph{Для пароля по этому словарю рекомендуется фраза не короче восьми случайно выбранных слов.}
  \TextParagraph{Electronic Frontier Foundation (EFF) рекомендует не менее шести слов из своего большого списка на 7776
    слов \hyperlink{source-2}{[2]}. При выборе восьми слов из списка на 1296 слов возможных фраз больше, чем при выборе
    шести слов из списка на 7776 слов. Поэтому восемь слов позволяют получить не меньшую стойкость к подбору.}
  \TextParagraph{Для дополнительной защиты количество слов можно увеличить.}

  \section{Состав словаря}
  \TextParagraph{Словарь позволяет запоминать пароль через русские слова и выбирать удобный способ его представления:}

  \begin{List}
    \BulletItem{Русские слова служат основой для запоминания парольной фразы.}
    \BulletItem{Транслитерация позволяет записать фразу строчными латинскими буквами. Не все системы поддерживают
      кириллицу, а правила ее обработки могут различаться между системами. Запись латиницей помогает избежать проблем с
      совместимостью.}
    \BulletItem{Трехбуквенные сокращения удобны для паролей, которые приходится часто вводить вручную. Все сокращения в
      словаре уникальны, поэтому их использование не уменьшает число возможных паролей. Для восстановления пароля
      достаточно помнить слова в правильном порядке и издание словаря.}
    \BulletItem{Трехзначные числовые коды позволяют создавать и запоминать длинные числа, например PIN-коды. Коды от 000
      до 999 присвоены 1000 словам.}
  \end{List}

  \TextParagraph{Словарь не является секретом. В тайне должна оставаться сама фраза. Транслитерация и сокращения служат
    для представления фразы и не являются шифрованием.}

  \subsection{Издание и редакция}
  \TextParagraph{На обложке указаны издание и редакция словаря. В новом издании могут измениться состав слов, их
    соответствие комбинациям кубиков, транслитерации, сокращения и числовые коды. Для восстановления пароля потребуется
    то же издание, поэтому его номер следует запомнить, а копию словаря — сохранить.}
  \TextParagraph{Редакция меняет только оформление и пояснения. Для восстановления пароля подходит любая редакция того
    же издания.}

  \section{Транслитерация}
  \TextParagraph{Для каждого слова в словаре закреплен один вариант транслитерации, чтобы одна и та же фраза всегда
    давала одинаковую запись латиницей. Этот вариант приведен рядом со словом. Транслитерация построена по общей
    системе, правила которой могут помочь запомнить написание слов без постоянного обращения к словарю.}

  \begin{TransliterationTable}
    \begin{TransliterationBlock}
      \TransliterationRow{а}{a}
      \TransliterationRow{б}{b}
      \TransliterationRow{в}{v}
      \TransliterationRow{г}{g}
      \TransliterationRow{д}{d}
      \TransliterationRow{е}{e \emph{или} ye}
      \TransliterationRow{ё}{yo}
      \TransliterationRow{ж}{zh}
      \TransliterationRow{з}{z}
      \TransliterationRow{и}{i}
      \TransliterationRow{й}{y}
    \end{TransliterationBlock}\hspace{\TransliterationBlockGap}%
    \begin{TransliterationBlock}
      \TransliterationRow{к}{k}
      \TransliterationRow{л}{l}
      \TransliterationRow{м}{m}
      \TransliterationRow{н}{n}
      \TransliterationRow{о}{o}
      \TransliterationRow{п}{p}
      \TransliterationRow{р}{r}
      \TransliterationRow{с}{s}
      \TransliterationRow{т}{t}
      \TransliterationRow{у}{u}
      \TransliterationRow{ф}{f}
    \end{TransliterationBlock}\hspace{\TransliterationBlockGap}%
    \begin{TransliterationBlock}
      \TransliterationRow{х}{h}
      \TransliterationRow{ц}{ts}
      \TransliterationRow{ч}{ch}
      \TransliterationRow{ш}{sh}
      \TransliterationRow{щ}{shch}
      \TransliterationRow{ъ}{\emph{пропускается}}
      \TransliterationRow{ы}{y}
      \TransliterationRow{ь}{\emph{пропускается}}
      \TransliterationRow{э}{e}
      \TransliterationRow{ю}{yu}
      \TransliterationRow{я}{ya}
    \end{TransliterationBlock}
  \end{TransliterationTable}

  \TextParagraph{Буква «е» обычно записывается как «e», а сразу после «ь» или «ъ» — как «ye». Например, слову «варенье»
    соответствует запись «varenye».}

  \section{Трехбуквенные сокращения}
  \TextParagraph{Сокращение состоит из трех строчных латинских букв транслитерации, расположенных в исходном порядке.}
  \TextParagraph{Для большинства слов используется следующее правило:}

  \begin{List}
    \NumberedItem{1}{Первой берется первая буква транслитерации.}
    \NumberedItem{2}{Второй берется ближайшая следующая согласная, если она не последняя буква транслитерации. Если она
      последняя или согласных больше нет, берется вторая буква транслитерации.}
    \NumberedItem{3}{Третьей берется ближайшая согласная после выбранной второй буквы. Если согласных дальше нет,
      берется буква, стоящая сразу после выбранной второй.}
  \end{List}

  \TextParagraph{При этом гласными считаются «a», «e», «i», «o», «u», а также «y», если она передает русскую «ы». В
    остальных случаях «y» считается согласной. Сочетания «zh», «ch», «shch», «yo» и другие обрабатываются побуквенно, а
    не как единый знак.}
  \TextParagraph{Например, слово «молоко» записывается как «moloko» и сокращается до «mlk», а «абрикос» записывается как
    «abrikos» и сокращается до «abr». Для слова «море» сокращение «mre» включает конечную гласную «e», поскольку после
    «r» согласных больше нет. В слове «боб», записанном как «bob», следующая согласная стоит в самом конце. Поэтому
    второй берется «o», а третьей — конечная «b».}
  \TextParagraph{Однако единое правило иногда дает одинаковые сокращения для разных слов. Чтобы все сокращения
    оставались уникальными, для части слов выбраны другие сочетания букв. Например, «автомобиль» имеет сокращение «avt»,
    а «автобус» — «avs».}
  \TextParagraph{Общее правило помогает запоминанию, но для точного восстановления пароля следует сверяться с
    сокращением, указанным в словаре.}

  \section{Создание парольной фразы}
  \subsection{Подготовка к броскам}
  \TextParagraph{Проще всего использовать четыре обычных шестигранных кубика. До первого броска определите, в каком
    порядке будете читать значения.}

  \begin{List}
    \BulletItem{Если кубики одинаковые, можно читать их слева направо относительно своего места за столом.}
    \BulletItem{Если кубики различаются, например по цвету, заранее назначьте каждому постоянную позицию: первую,
      вторую, третью или четвертую.}
  \end{List}

  \TextParagraph{Порядок чтения остается неизменным для всех бросков и не зависит от величины выпавших значений.}
  \TextParagraph{Удобно встряхивать кубики в стаканчике и бросать на ровную поверхность в лотке с невысокими бортиками.
    Кубикам должно хватать места, чтобы свободно перекатываться и останавливаться плашмя. Вместо стаканчика можно
    использовать башню для кубиков.}

  \subsection{Выбор слов}
  \begin{List}
    \NumberedItem{1}{Бросьте четыре кубика и прочитайте результат в заранее установленном порядке. Если кубик встал на
      ребро, вылетел из лотка или порядок чтения кубиков неясен, перебросьте все четыре кубика.}
    \NumberedItem{2}{Найдите выпавшую комбинацию в столбце «Кубики». Соответствующее слово станет следующим словом
      фразы.}
    \NumberedItem{3}{Запомните или запишите слово и повторите бросок для следующего.}
    \NumberedItem{4}{Продолжайте, пока не получите нужное количество слов.}
  \end{List}

  \TextParagraph{Сохраняйте все выпавшие слова в исходном порядке, включая повторяющиеся. Замена слов и дополнительные
    броски ради более удобного или приятного варианта нарушают случайность выбора и снижают стойкость фразы.}

  % Editorial page break keeps this short subsection together on its own page.
  \newpage
  \subsection{Запоминание фразы}
  \TextParagraph{Свяжите выпавшие слова в историю или мысленную сцену, которая помогает вспомнить их в исходном порядке.
    История может быть нелепой и необычной.}
  \TextParagraph{На время создания и запоминания можно записать слова на бумаге. Держите запись в безопасном месте.
    Когда убедитесь, что уверенно помните фразу, временную записку рекомендуется уничтожить так, чтобы запись нельзя
    было восстановить.}

  \section{Создание пароля}
  \TextParagraph{Сначала получите фразу не короче восьми слов по инструкции выше. Затем выберите один из двух способов
    ее представления в качестве пароля.}

  \subsection{Пароль из сокращений}
  \TextParagraph{Для каждого слова возьмите сокращение из словаря и соедините сокращения подряд, без пробелов и других
    разделителей. Сохраняйте порядок слов. Используйте строчные латинские буквы.}
  \TextParagraph{У каждого сокращения ровно три буквы, поэтому границы слов определяются однозначно даже без пробелов.
    Пароль из восьми слов будет содержать 24 буквы.}

  \subsection{Пароль из полных транслитераций}
  \TextParagraph{Для каждого слова возьмите готовую транслитерацию из словаря. Соедините записи в том же порядке,
    разделяя их одним пробелом. Если система не принимает пробелы, используйте один и тот же разрешенный небуквенный
    символ, например дефис, между всеми словами.}
  \TextParagraph{Разделитель нужен для однозначного обозначения границ слов. Например, фразы «пол экран» и «поле кран»
    без разделителя дают одинаковую запись «polekran». Запомните выбранный разделитель вместе с фразой.}

  \subsection{Требования к составу пароля}
  \TextParagraph{Если система требует заглавных букв, цифр или специальных символов, можно заранее выбрать постоянное
    правило дополнения пароля. Например, первую букву делать заглавной, а в конце добавлять одну и ту же цифру и
    специальный символ, допустимые в системе.}
  \TextParagraph{Такое дополнение служит только для выполнения формальных требований. Стойкость пароля следует считать
    равной стойкости исходной случайной фразы, не рассматривая правило дополнения как дополнительный секрет.}

  \section{Создание PIN-кода}
  \TextParagraph{Для PIN-кода слова служат подсказкой для запоминания числа.}

  \begin{List}
    \NumberedItem{1}{Определите, сколько слов потребуется. Каждое слово дает три цифры. Для PIN-кода из четырех или
      шести цифр достаточно двух слов.}
    \NumberedItem{2}{Выбирайте слова бросками четырех кубиков, как описано выше, но используйте только слова с
      заполненным столбцом «Код». Если у слова нет числового кода, пропустите его и перебросьте все четыре кубика.}
    \NumberedItem{3}{Для каждого выбранного слова возьмите трехзначный код. Соедините коды подряд в той же
      последовательности, что и слова во фразе.}
    \NumberedItem{4}{Если цифр получилось больше, чем нужно, отбросьте лишние справа.}
    \NumberedItem{5}{Проверьте получившийся PIN-код по правилам следующего раздела. Если он не подходит, начните
      создание заново.}
  \end{List}

  \TextParagraph{Все начальные нули в кодах и готовом PIN-коде сохраняются.}
  \TextParagraph{Для восстановления достаточно помнить слова в правильном порядке, длину PIN-кода и издание словаря. Для
    PIN-кода из четырех цифр вместо двух слов можно запомнить первое слово и четвертую цифру полученного числа: первые
    три цифры восстанавливаются по коду слова, а последняя — по памяти.}

  \subsection{Проверка PIN-кода}
  \TextParagraph{В редких случаях сгенерированный PIN-код может совпасть с одной из комбинаций, часто используемых при
    подборе. В качестве дополнительной предосторожности не используйте PIN-коды, которые:}

  \begin{List}
    \BulletItem{состоят из одинаковых цифр, например 000000}
    \BulletItem{состоят из цифр, идущих подряд по возрастанию или убыванию с шагом один, например 012345 или 765432}
    \BulletItem{состоят из повторяющихся блоков, например 121212 или 123123}
    \BulletItem{состоят из последовательных пар одинаковых цифр, например 112233}
    \BulletItem{одинаково читаются слева направо и справа налево, например 123321}
    \BulletItem{совпадают с вашей датой рождения, записанной в одном из распространенных форматов}
  \end{List}

  \TextParagraph{Если получился такой PIN-код, начните весь процесс заново, с новых бросков для первого слова. Не
    исправляйте отдельные цифры вручную.}

  \section{Источники}
  \begin{List}
    % Keep URLs freely wrapped instead of tightening the right edge as in prose.
    \raggedright
    \NumberedItem{1}{\ReferenceTarget{source-1}Reinhold, A. G. The Diceware Passphrase Home Page.
      \url{https://theworld.com/~reinhold/diceware.html} (дата обращения: 16.10.2026).}
    \NumberedItem{2}{\ReferenceTarget{source-2}Electronic Frontier Foundation. EFF Dice-Generated Passphrases.
      \url{https://www.eff.org/dice} (дата обращения: 16.10.2026).}
  \end{List}

  \TableInitializeVerticalAlignment
  \TableStartOnLeftPage
  \phantomsection
  \addcontentsline{toc}{section}{Словарь}
  \begin{DictionaryTable}
    \TableHeader{Кубики}{Слово}{Транслитерация}{Сокр.}{Код}
    \input{dictionary-rows.tex}
  \end{DictionaryTable}
\end{document}
